\scanpage{95}

que $N(x)$ induit sur $R^n$ la fonction $N_0(x)$, et qu'on a $N(\lambda x)=|\lambda|N(x)$ et $N(x)=0\Rightarrow x=0$, il faut seulement prouver $N(x+y)\leq N(x)+N(y)$ (\unsure{en\textsuperscript{t}|est|seule} valable $=$ pour $x,y\in c_0$). Pour ceci, on \deleted{\unsure{regarde $(x_i)$ et $(y_i)$ comme les}}\inserted{remarque que si} $u_x$ est l'opération de multiplication par $x$ dans $\ell^2$ (qui est un opérateur cpt pour $x\in c_0$), on a \deleted{$u_{x+y}=u_x$} $\rho_i(u_x)=\rho_i(x)$, d'autre part \deleted{\unsure{que}} $u_{x+y}=u_x+u_y$, et l'inégalité à prouver résulte alors du N\textsuperscript{o}~4, th.~4, corollaire. En $=$ \deleted{\unsure{temps}}, les normes \inserted{de Schatten} ainsi associées à des normes \deleted{\unsure{de}} de Schatten $N_0$ sur $R^{n}$ \deleted{donnent} démontrent qu'il existe une infinité \inserted{continue} de normes de Schatten sur $R^{(\underline N)}$ équivalentes à la norme induite par $c_0$.

Soit $N$ \deleted{est} une norme de Schatten sur \deleted{\unsure{$R^n$} \unsure{finies}} $R^{(\underline N)}$. \deleted{Soit $\ell^N$ \unsure{l'espace des}}\inserted{Pour \deleted{toute}} suite \deleted{\unsure{réelle}} complexe $x=(x_i)$ \deleted{telle que \unsure{la suit} \unsure{on pose} soit $\unsure{N(x\ldots x)|N(x)}$}

\[
\deleted{N(|x_1|,\ldots |x_n|,0,0,\ldots)\quad\text{croît avec }n,\quad\text{on pose.}}
\]

Soit $|x|_n=(|x_1|,\ldots |x_n|,0,\ldots 0)$. \deleted{On a \unsure{aussitôt}}\inserted{Si $n$ varie} \deleted{alors $|x+y|_n\leq |x|_n+|y|_n$} $|x|_n$ est une suite croissante d'éléments positifs de $R^{(N)}$, posons

\[
N(x)=\lim_n N(|x|_n)\tag*{(4)}
\]

c'est un nb $\geq 0$ fini ou infini. La valeur \deleted{ainsi \unsure{désignée}}\inserted{ainsi} de $N(x)$ coïncide avec \deleted{\unsure{celle|$N(x)$}} \inserted{quand $x\in R^{(N)}$.} On a

\[
N(\lambda x)=|\lambda|N(x),\qquad N(\unsure{u|x}+y)\leq N(\unsure{u|x})+N(y)\tag*{(5)}
\]

\deleted{si $N(x)<+\infty$} \inserted{et dans} la première formule, \deleted{\unsure{il est à dire que son deuxième mb a un sens, i.e. $|\lambda|$}} \inserted{il faut prendre au deuxième} mb la valeur \uline{zéro} s'il \deleted{\unsure{est sous}}\inserted{se présente sous la} forme indéterminée $0\times\infty$. La deuxième formule résulte de ce que \deleted{\unsure{$N(x$}} $|x+y|_n\leq |x|_n+|y|_n$, d'où $N(|x+y|_n)\leq N(|x|_n+|y|_n)\leq N(|x|_n)+N(|y|_n)$.

\margin{\unsure{?}}

\deleted{Soit} Soit $\ell^N$ l'espace des suites complexes $x$ telles que $N(x)<\infty$, \unsure{muni évidemment des notions|espace vectoriel au moins|muni de} \deleted{\unsure{celles}} \unsure{des suites complexes nulles|de suite complexe nulle}, \unsure{partie|d'opération} \deleted{\unsure{de mult}} d'un \unsure{espace vect.|certain rang}. \unsure{Par|Sur|Pour} $(5)$, on voit que $N(x)$ est une \uline{norme} sur $\ell^N$.

\deleted{\unsure{De plus|D'ailleurs}, $\ell^N$ est un espace de \illegible{2} \unsure{positivité}. D'où \unsure{les inj} $R^{(N)}\in\ell^N$, et $N(x)$ induit \unsure{par hypothèse}, \unsure{donc} ($=$ \illegible{1} $\unsure{\alpha^{(i)}}$), est \illegible{2} \unsure{le candidat} dans $\ell^N$, \unsure{elles coïncident donc forcément}}
% archive sheet number: 96
